\section{Introduction}
\label{Introduction}

Single vehicle autonomy, which relies on camera and LiDAR observations, is fundamentally limited by local sensing, including occlusions, a finite range, and partial observability in dense urban traffic. For example, a pedestrian stepping out from behind a parked truck may be invisible to the ego vehicle but clearly visible to a nearby roadside unit. Connected autonomous vehicles (CAVs) address such limitations by sharing information through Vehicle-to-Everything (V2X) communication~\cite{han2023collaborative, huang2023v2xsurvey}. Collaborative perception fuses multi-agent observations to produce richer, more complete representations of the driving environment~\cite{teufel2024collective, wang2020v2vnet, HuEtAl2022}.

V2Xverse~\cite{liu2024v2xverse} advances this line of work by framing collaboration as a full-stack closed-loop driving problem rather than a perception-only benchmark. Its CoDriving pipeline uses collaborative features to improve waypoint prediction and route-level behavior in CARLA. However, stronger perception and planning do not guaranty safe outcomes in themselves. The low-level controller that converts planned waypoints into throttle and brake commands has no explicit model of surrounding actors: it tracks the desired speed faithfully, regardless of whether a vehicle is cutting in or a cyclist is approaching from the side. Safety failures in this stack are therefore a control-level problem, not a perception-level one: collisions occur even when the upstream perception correctly detects all relevant actors.

Control Barrier Functions (CBFs) offer a mechanism to enforce safety at the control level by restricting actions to remain within a forward invariant safe set~\cite{ames2017cbf, ames_control_2019}. A CBF quadratic program (QP) minimally modifies the nominal command to satisfy the barrier constraints. A central parameter is the barrier gain $\gamma$, which determines how strongly the filter enforces the safety boundary. Its effect is state-dependent: near the safety boundary or under unfavorable closing dynamics, a larger $\gamma$ triggers earlier intervention, whereas in safer states the same gain has little effect. A fixed $\gamma$ therefore imposes a single operating point that must simultaneously handle cruise, dense vehicle interactions, pedestrian crossings, and bicycle conflicts. Conservative settings protect safety but reduce throughput; permissive settings preserve progress but under-react to rapidly evolving hazards. This motivates learning a context-dependent gain policy that adapts barrier enforcement online.

Retraining the upstream planner with a safety-aware objective, given its access to full V2X perception, is a natural alternative to a separate safety layer. We deliberately choose a frozen-planner adapter instead, for three reasons. First, parameter efficiency: our CBF-gain policy has on the order of $10^4$ trainable parameters, roughly three orders of magnitude fewer than the collaborative perception-planning stack it sits downstream of. Second, risk isolation: retraining the planner risks degrading an already-validated perception or planning behavior on its primary task, whereas a frozen-planner adapter cannot regress planning quality by construction. Third, portability: the safety layer attaches to any upstream planner's output trajectory without requiring access to that planner's training pipeline, weights, or architecture, a property a retrained, safety-aware planner would not have.

We address this by integrating adaptive safety modulation into V2Xverse, employing constrained reinforcement learning to learn state-dependent CBF gains. Building on recent work that learns adaptive barrier gains in simplified multi-agent settings~\cite{berducci2024las}, we keep collaborative perception and planning frozen and train only a compact policy ($\sim$10k parameters) that outputs class-specific CBF gains (one each for vehicles, pedestrians, and bicycles) at every decision step. The learned gains parameterize a CBF-QP filter~\cite{ames2017cbf} that adjusts longitudinal control while preserving the nominal steering path. Training uses a constrained PPO variant~\cite{schulman2017ppo} with a PID-based Lagrangian multiplier~\cite{stooke2020pid} to satisfy a cumulative safety cost budget. Because the upstream stack is untouched, the training is sample-efficient, and any change in safety or efficiency is attributable solely to the adaptive control layer. Unlike end-to-end RL approaches that learn perception, planning, and control jointly, our design is interpretable (the output is a small set of physically meaningful gains), composable (it wraps any nominal controller), and lightweight to train.

The contributions of this paper are:
\begin{itemize}
    \item We present an adaptive CBF safety layer integrated into the closed-loop V2Xverse CoDriving stack, adding learned safety modulation without retraining perception or planning.
    \item We formulate a class-aware adaptive gamma policy trained with constrained PPO and a PID-based Lagrangian~\cite{stooke2020pid} that learns state-dependent CBF gains from compact traffic risk features.
    \item We evaluate on CARLA Town05 under multi-scenario routing, comparing adaptive CBF against fixed CBF and unfiltered baselines under identical perception and planning.
    \item We show that adaptive CBF modulation reduces collision rates while maintaining route completion, demonstrating that safety can be improved by only learning gain scheduling.
\end{itemize}

The remainder of this paper is organized as follows. Section~\ref{Related_work} reviews related work. Section~\ref{methods} presents the adaptive CBF method. Section~\ref{experiments} describes the experimental setup. Section~\ref{results} reports results. Section~\ref{Conclusion} concludes.
